% Part: first-order-logic
% Chapter: syntax-and-semantics
% Section: main-operator

\documentclass[../../../include/open-logic-section]{subfiles}

\begin{document}

% SEGMENT OLP-0154-S01
\olfileid{fol}{syn}{mai}

\olsection{\printtoken{S}{main operator} ஒரு சூத்திரத்தில்}

\begin{explain}
!!a{formula}~$!A$ ஐ உருவாக்கப் பயன்படுத்திய கடைசி இயக்கியைப் பற்றி பேசுவது
பல நேரங்களில் பயனுள்ளதாக இருக்கும். இந்த இயக்கி~$!A$ இன் \emph{முதன்மை
இயக்கி} எனப்படுகிறது. உள்ளுணர்வாக, இது $!A$ இன் ``வெளிப்புற'' இயக்கி.
உதாரணமாக, $\lnot !A$ இன் முதன்மை இயக்கி $\lnot$; $(!A \lor !B)$ இன்
முதன்மை இயக்கி $\lor$; இவ்வாறே.
\end{explain}

\begin{defn}[!!^{main operator}]
\ollabel{def:main-op}
!!a{formula}~$!A$ இன் \emph{!!{main operator}} பின்வருமாறு வரையறுக்கப்படுகிறது:
\begin{enumerate}
\item \indcase*{!A}{!A}{$\indfrm$ க்கு !!{main operator} இல்லை.}

\tagitem{prvNot}{\indcase{!A}{\lnot !B}{$\indfrm$ இன் !!{main operator}
  $\lnot$.}}{}

\tagitem{prvAnd}{\indcase{!A}{(!B \land !C)}{$\indfrm$ இன் !!{main operator}
  $\land$.}}{}

\tagitem{prvOr}{\indcase{!A}{(!B \lor !C)}{$\indfrm$ இன் !!{main operator}
  $\lor$.}}{}

\tagitem{prvIf}{\indcase{!A}{(!B \lif !C)}{$\indfrm$ இன் !!{main operator}
  $\lif$.}}{}

\tagitem{prvIff}{\indcase{!A}{(!B \liff !C)}{$\indfrm$ இன் !!{main operator}
  $\liff$.}}{}

\tagitem{prvAll}{\indcase{!A}{\lforall[x][!B]}{$\indfrm$ இன் !!{main operator}
  $\lforall$.}}{}

\tagitem{prvEx}{\indcase{!A}{\lexists[x][!B]}{$\indfrm$ இன் !!{main operator}
  $\lexists$.}}{}
\end{enumerate}
\end{defn}

ஒவ்வொரு நிலையிலும், சூத்திரத்தில் சுட்டிக்காட்டப்பட்ட !!{main operator} இன்
குறிப்பிட்ட \emph{நிகழ்வையே} குறிக்கிறோம். உதாரணமாக,
$((!D \lif !E) \lif (!E \lif !D))$ என்பது $(!B \lif !C)$ வடிவில் உள்ளது;
இங்கு $!B$ என்பது $(!D \lif !E)$, $!C$ என்பது $(!E \lif !D)$. ஆகவே
$\lif$ இன் இரண்டாவது நிகழ்வே !!{main operator}.

\begin{explain}
இது அணுவல்லாத !!{formula}s அனைத்தையும் அவற்றின் !!{main operator}
நிகழ்வுகளுக்கு மாற்றும் ஒரு சார்பின் \emph{மறுநிகழ்வு} வரையறை. !!{formula}s
தொகுத்தறிதல் வழியில் வரையறுக்கப்பட்ட விதத்தால், ஒவ்வொரு !!{formula}~$!A$ வும்
\olref{def:main-op} இன் ஏதாவது ஒரு நிலையில் அடங்கும். எனவே அணுவல்லாத
ஒவ்வொரு !!{formula}~$!A$ க்கும் !!a{main
  operator} இருப்பது உறுதி.
ஒவ்வொரு !!{formula} வும் இந்நிபந்தனைகளில் ஒன்றை மட்டுமே பூர்த்தி செய்கிறது;
மேலும் ஒவ்வொரு நிலையில் $!A$ உருவாக்கப்பட்ட சிறிய !!{formula}s தனித்துவமாகத்
தீர்மானிக்கப்படுகின்றன. ஆகவே $!A$ இன் !!{main
  operator} நிகழ்வு தனித்துவமானது;
இதனால் ஒரு சார்பை வரையறுத்துள்ளோம்.
\end{explain}

!!{formula}s இன் !!{main operator} எந்தக் குறி என்பதைப் பொறுத்து,
\olref{tab:main-op} இல் உள்ள பெயர்களால் அழைக்கிறோம்.\iftag{defNot,defOr,defAnd,defIf,defIff,defTrue,defFalse,defEx,defAll}
{ ஆனால் வரையறுக்கப்பட்ட இயக்கிகள் அதிகாரப்பூர்வமாக !!{formula}s இல்
தோன்றுவதில்லை என்பதை நினைவில் கொள்ளுங்கள். அவை சுருக்கங்கள் மட்டுமே; எனவே
அதிகாரப்பூர்வமாக ஒரு சூத்திரத்தின் முதன்மை இயக்கியாக இருக்க முடியாது.
அனைத்து !!{formula}s பற்றிய நிறுவல்களில் அவற்றைத் தனித்தனியாகக் கருதத்
தேவையில்லை.}

\begin{table}[!h]
\centering
\begin{tabular}{c | c | c}
!!^{main operator} & !!{formula} இன் வகை & எடுத்துக்காட்டு\\
\hline
none & atomic (!!{formula}) &
\iftag{prvFalse}{$\lfalse$,}{}
\iftag{prvTrue}{$\ltrue$,}{}
$\Atom{R}{t_1, \dots, t_n}$,
$\eq[t_1][t_2]$\\
$\lnot$ & மறுப்பு & $\lnot !A$ \\
$\land$ & இணையல் & $(!A \land !B$) \\
$\lor$ & பிரிப்பிணைவு & $(!A \lor !B$) \\
$\lif$ & !!{conditional} & $(!A \lif !B$) \\
$\liff$ & !!{biconditional} & $(!A \liff !B)$ \\
$\lforall[][]$ & universal (!!{formula})& $\lforall[x][!A]$ \\
$\lexists[][]$ & existential (!!{formula})& $\lexists[x][!A]$
\end{tabular}
\caption{முதன்மை இயக்கியும் !!{formula}s இன் பெயர்களும்}
\ollabel{tab:main-op}
\end{table}

\end{document}
